% BOOK-FILE-ID: f-79edd838da43
% section: おまけ：定数数列

\section{定数型}\label{節：定数数列}
  \booknavref{解法：階比型の1次式の積}{1次式の積の場合}の中で,
  \begin{align*}
    g(n+1)a_{n+1}&=g(n)a_n\\
    &=g(n-1)a_{n-1}\\
    &=g(n-2)a_{n-2}\\
    &\qquad \vdots\\
    &=g(1)a_1\\
    &\therefore \quad a_n=\frac{g(1)a_1}{g(n)}
  \end{align*}
  のような式変形を扱った.
  ここで,$b_n=g(n)a_n$と置くと,
  \[
    b_n=b_{n-1}=b_{n-2}=\cdots=b_1
  \]
  のような式に書き換えることができる.
  このように,元の数列を調べるために新しく置いた数列を\textbf{\bookterm{auxiliary}{補助数列}}という.
  この$b_n$のように,常に一定の値を取り続ける数列を定数数列という.
  直前の式では,常に$b_1$の値を取っている.

  \begin{bookpoint}
    \textbf{POINT}\par
    $b_{n+1}=b_n$となる補助数列を作れば,すべての項は$b_1$に等しい.
    初項でその値を決め,元の数列へ戻す.割り算で戻すときは,分母が0でないことを確かめる.
  \end{bookpoint}
  
